\scanpage{32}
% L1–5
Nous nous intéressons notamment au cas où les $D(A)$ sont finis, et vides sauf un nb fini d'entre eux, i.e. $\exists N\in\mathbb N$ tel que $A\not\subset[0,N]\Longrightarrow D(A)=\varnothing$. Il revient au $=$ que $\Theta(D)$ est fini.
% L6–8
La réalisation topologique \inserted{$|\Theta(D)|$} de $\Theta(D)$ \deleted{\unsure{noté|notée}} est aussi notée $|D|$ et s'appelle réalisation topologique de $D$.
% L8–11
C'est un espace à triangulation finie, \deleted{\unsure{les}} et $\Theta(D)$ en tant qu'ens ordonné se \unsure{rélise|réalise} comme ensemble de parties de $|D|$.
% L12–14
On va donner une deuxième description de $|D|$, en définissant certaines opérations simpliciales entre les $\Theta_n(\unsure{\mathbb D|D})$.
% L15–17
Posons $\Delta_n=[0,n]\subset\mathbb N$, soit
\[
\sigma\colon\Delta_n\hookrightarrow\Delta_m
\]
une application strictm\textsuperscript{t} croissante. Si
% L18–20
$A\in\unsure{\mathscr P|P|\mathfrak P}_{\unsure{f\,m+1|f_{m+1}}}(\mathbb N)$, on a une unique bij. croissante
\[
\Delta_m\overset{u_A}{\underset{\longrightarrow}{\sim}} A
\]
% L20–22
d'où une partie finie de \deleted{\unsure{card $m+1$|card.}} \inserted{cardinal $n+1$}, $B=u_A\sigma(\Delta_n)$, de $A$, et on considère
\[
\sigma^*\colon\{\;D(A)\longrightarrow D(B)\hookrightarrow\Theta_n(D)
\qquad (B\mathbin{:}u_A\sigma(\Delta_{\unsure{n}})
\]
% archive sheet number: 31
